\section{Blackwell's general discounted model}
\label{sec:annex-p17-blackwell-general-discounted}

\begin{annexitem}{annex:p17.blackwell.factorization}{thm:p17-blackwell-factorization}{The exact kernel and product interfaces}
\begin{proof}
Each coordinate space has a countable topological base.
Finite-coordinate rectangles from those bases give a countable
base for a finite or countable product. Every product-open
set is therefore a countable union of measurable cylinders.
Conversely the coordinate projections are continuous.
This proves equality of the product and Borel sigma algebras.
A countable product of Borel subsets of Polish ambient spaces
is the intersection of their coordinate inverse images in
the Polish ambient product, hence is itself standard Borel.
For completeness of the ambient product, use complete coordinate
metrics \(d_i\) and the metric
\(\sum_{i\geq1}2^{-i}\min(1,d_i(x_i,y_i))\).
A Cauchy sequence is Cauchy in each coordinate, hence has
coordinate limits. A finite-head estimate followed by a
uniformly small geometric tail proves convergence in this
product metric. Its topology is the product topology, and the
countable base already exhibited gives separability.

On \((X\times Y,m)\), take the coordinate random variables.
The conditional law of the second given the first, constructed in
\cref{thm:p15-value-conditional-law}, is a kernel \(k\).
That theorem applies because the target \(Y\) has the Borel real
coding of \cref{thm:p15-metric-borel-representation}.
Its rectangle identity is
\(m(C\times D)=\int_C k(D\mid x)p(dx)\), where \(p\) is
the first marginal. The joint-measure theorem
\cref{thm:p15-kernel-joint-measure} gives equality on the whole
product sigma algebra and the formula for every bounded joint
integrand. Its proof by measurable sections, simple approximation
and monotone convergence also proves that integration preserves
measurability of a parameter in the integrand. The uniqueness
part of the conditional-law theorem uses one countable generating
class for \(Y\), so its exceptional union is one \(p\)-null set
on which equality as probability measures can fail.
The first marginal is forced by \(m(C\times Y)\).

For an additional parameter \(z\), start from the given initial
probability kernel and form successive finite laws by integrating
the later kernels in their prescribed order. Kernel
integration proves by induction that the probability of each
finite-dimensional Borel event is a measurable function of \(z\).
For clarity, in the induction the integrand for an event \(E\)
is its indicator on the whole prefix, not a presumption that
every event is a rectangle. The kernel-joint construction first
handles rectangles and extends to all such indicators.
Each finite law is a probability. Integrating out the last
coordinate, whose kernel has total mass one, proves consistency.

For fixed \(z\), \cref{thm:p15-coded-space-extension} now
gives a unique law on the countable coordinate product. Its
hypotheses hold for every coordinate space: each is standard
Borel, its code has Borel image, and the preceding finite laws
are consistent. Let \(\mathcal D\) be the class of product events
whose probabilities under these laws are measurable in \(z\).
It contains all finite cylinders. It contains the whole space,
is closed under complements by subtraction from one, and under
countable disjoint unions by the pointwise sum of probabilities.
The cylinder pi-system generates the product sigma algebra;
the pi-lambda theorem therefore makes \(\mathcal D\) the whole
sigma algebra. Thus the separately constructed laws form one
kernel. This also proves uniqueness of that kernel from its
finite distributions, at every parameter value.

Apply this to the alternating action and state kernels of the
policy, with \(z=s_1\), or with \(z=(s_1,\theta)\) for any
additional policy parameter. The reward series is a pointwise
limit of Borel functions and converges uniformly and absolutely:
its tail after \(N\) is bounded by
\(R\sum_{n>N}\beta^{n-1}=R\beta^N/(1-\beta)\).
At \(\beta=0\) only the first reward remains; for \(N=0\)
use the whole-return bound. Kernel integration and uniform
convergence give the bounded Borel value. Two policies agreeing
through day \(N\) have the same first-\(N\) reward expectation;
their two tails give the factor \(2\). These arguments use no
conditioning on a null history.
\end{proof}
\end{annexitem}

\begin{annexitem}{annex:p17.blackwell.positive.selection}{thm:p17-blackwell-positive-selection}{Complete topology refinement and selection in closed fibers}
\begin{proof}
\paragraph{The topology refinement used in the construction.}
We prove that a Borel subset of a Polish space admits a Polish
topology with the same Borel sets, and that any countable
collection of Borel sets can be made clopen by a finer such
topology. Here Polish means separable and completely metrizable.

First an open subset \(O\) of a complete metric space \(E\)
is completely metrizable. If its complement is nonempty, put
\(d_O(x,y)=d(x,y)+|1/d(x,E\setminus O)-1/d(y,E\setminus O)|\)
on \(O\). A \(d_O\)-Cauchy sequence converges in \(E\), and
the boundedness of its reciprocal distances keeps the limit
in \(O\); the displayed formula then gives convergence in
\(d_O\). It has the original topology because the reciprocal
distance is continuous on \(O\). If \(O=E\), retain \(d\).
A closed subspace is complete in \(d\).
All these subspaces are separable: a countable base for the
original separable metric space restricts to a countable base,
and choosing one point from each nonempty basic open set gives
a countable dense set.
Consequently a closed set and its open complement, given their
topological disjoint union, make a Polish topology on \(E\)
in which that closed set is clopen. A complete metric is obtained
by bounding each component metric by \(1\) and giving different
components distance \(2\). This topology is finer and has the
same Borel sets.

For countably many finer Polish topologies \(\tau_n\) on \(E\),
each with the original Borel sets, their common refinement is
Polish with the same Borel sets. Indeed include the original
topology as \(\tau_0\), and take the diagonal
\[
 \Delta=\{(x_0,x_1,\ldots)\in\prod_{n\geq0}(E,\tau_n):
                                      x_n=x_0\text{ for all }n\}.
\]
This set is closed: the identity maps into \((E,\tau_0)\)
are continuous, and equality is closed in a metric space.
The countable product is completely metrizable by
\(\sum_n2^{-n-1}\min(1,d_n(x_n,y_n))\).
Coordinatewise completeness, followed by a finite-head and
small-tail estimate, proves completeness for this metric.
Finite-coordinate products of countable bases give a countable
base and separability. The closed diagonal is therefore Polish.
Pull its topology back by the diagonal bijection.
It is the common refinement. Its countable basic opens are
finite intersections of original Borel sets, so its Borel
sigma algebra is exactly the original one.

Let \(\mathcal C\) consist of subsets of \(E\) that can be
made clopen by such a refinement. It contains every closed
set by the disjoint-union construction. It is closed under
complements. If \(B_n\in\mathcal C\), take the common refinement
of topologies making each \(B_n\) clopen. Their union is open
in that Polish topology; the disjoint-union construction applied
once more makes the union clopen. Hence \(\mathcal C\) is a
sigma algebra containing the closed sets, and contains every
Borel set. Apply a common refinement to any countable list
to make them all clopen simultaneously.
In particular a Borel subset \(B\) becomes a closed Polish
subspace in a refinement of its ambient space, with its
original trace Borel sets unchanged.

\paragraph{A Polish graph with closed fibers.}
Give \(X,Y\) Polish topologies with their specified Borel sets,
as just proved, and view \(D\) as a standard Borel graph space.
Give \(D\) a Polish topology with its Borel sets unchanged.
The projection \(p:D\to X\) is Borel. For a countable base
\((V_j)\) of \(X\), make all \(p^{-1}(V_j)\) clopen by the
countable refinement construction. Then \(p\) is continuous.
In particular each fiber \(F_x=p^{-1}\{x\}\) is closed in this
complete metric graph space.

Put \(d(x)=k(D_x\mid x)>0\). The kernel section theorem
makes \(d\) Borel. Define a probability kernel on \(D\) by
\[
       \nu_x(C)=\frac{1}{d(x)}
           \int_Y\mathbf1_{C}(x,y)\mathbf1_D(x,y)k(dy\mid x).
\]
A Borel subset of \(D\) is the trace of a Borel subset of
\(X\times Y\), so the formula is unambiguous and Borel in
\(x\); normalized restriction gives countable additivity and
total mass one. Moreover \(\nu_x(F_x)=1\).

\paragraph{A Borel limit in each fiber.}
Fix a compatible complete metric on \(D\), and enumerate all
open balls with centers in a countable dense set and positive
rational radii. This is a base. Choose \(O_1(x)\) to be the
first such ball of diameter at most \(1/2\) with
\(\nu_x(O_1(x))>0\). At step \(n+1\), choose the first ball
of diameter at most \(2^{-n-1}\), with closure contained in
\(O_n(x)\), and of positive \(\nu_x\)-measure.
Such a ball exists: the allowed balls cover the preceding
open set, and a countable cover of a positive-measure set
cannot have all members null. The initial existence follows
from a countable cover of all of \(D\).
At each step the chosen index is a Borel function of \(x\).
To see this inductively, restrict to each Borel set on which
the preceding finite indices are fixed; the admissible
containment and diameter conditions are then fixed facts
about basis balls, and positive measure is a Borel condition.
Taking the first admissible index is a countable Borel partition.

The centers \(c_n(x)\) are Borel maps with countable range.
Nested closures and shrinking diameters make them Cauchy;
completeness gives a limit \(z(x)\in D\), in every
\(\overline{O_n(x)}\). This limit map is Borel: for any nonempty
closed set \(F\), the real functions \(d(c_n(x),F)\) are Borel and
converge to \(d(z(x),F)\); its zero set is the inverse image
of \(F\). The empty set has empty inverse image.
Closed sets generate the Borel sigma algebra.
Each chosen open ball meets the closed fiber \(F_x\),
because it has positive measure under a probability concentrated
on \(F_x\). Choose for this fixed \(x\) a point of
\(F_x\cap O_n(x)\); its distance to \(z(x)\) tends to zero.
Closedness therefore gives \(z(x)\in F_x\).
Finally the second-coordinate map \(D\to Y\) is Borel for
the unchanged Borel sets. Its composite with \(z\) is the
required \(f\), and \((x,f(x))\in D\) for every \(x\).
No measurability is asserted for the auxiliary fiber points
used only to prove this last membership.
\end{proof}
\end{annexitem}

\begin{annexitem}{annex:p17.blackwell.selection}{thm:p17-blackwell-selection}{Mean improvement and the strict essential-supremum margin}
\begin{proof}
Let \(m(x)\) be the essential supremum of \(v(x,\cdot)\)
under \(k(\cdot\mid x)\). It is a finite Borel function:
\[
 m(x)=\sup\{t\in\mathbb Q:
                     k(\{y:v(x,y)>t\}\mid x)>0\}.
\]
All the threshold probabilities are Borel by the joint
kernel-integration theorem. The set inside the supremum
contains every sufficiently negative rational and is bounded
above by \(\|v\|_\infty\). For fixed \(x\), the definition
implies \(v(x,\cdot)\leq m(x)\) almost surely, by taking the
countable union of the null rational upper tails.
Also every threshold strictly below \(m(x)\) has positive
upper-tail probability. Let \(a(x)=(kv)(x)\leq m(x)\).

Use the following Borel set:
\[
 D=\{(x,y):a(x)<m(x),\
              v(x,y)\geq a(x),\ v(x,y)>m(x)-\eta\}
       \ \cup\
   \{(x,y):a(x)=m(x),\ v(x,y)=m(x)\}.
\]
If \(a(x)<m(x)\), the threshold
\(\max(a(x),m(x)-\eta)\) is strictly below \(m(x)\).
Its strict upper tail is contained in \(D_x\) and has
positive probability. If \(a(x)=m(x)\), the almost surely
nonnegative random variable \(m(x)-v(x,\cdot)\) has integral zero,
so it is zero almost surely. Thus \(D_x\) has probability
one in this case. The positive-section lemma supplies a
Borel selection \(f\).
In either case \(v(x,f(x))\geq a(x)\) and
\(v(x,f(x))>m(x)-\eta\). Hence
\(v(x,f(x))+\eta>m(x)\), and the set with
\(v(x,y)\geq v(x,f(x))+\eta\) is null.
The strict margin is why the statement includes the
equality threshold in its null event.

For a Borel \(f\), \(\delta_{f(x)}(C)=\mathbf1_C(f(x))\)
is a degenerate kernel. Conversely, if a kernel is a
point mass at a unique point \(f(x)\) for every \(x\),
then for every Borel \(C\),
\(f^{-1}(C)=\{x:k(C\mid x)=1\}\) is Borel. Thus \(f\)
is Borel. Uniqueness of the point follows because distinct
points in a metric space are separated by Borel sets.
\end{proof}
\end{annexitem}

\begin{annexitem}{annex:p17.blackwell.prior.optimal}{thm:p17-blackwell-prior-optimal}{The integrated supremum and countable initial-state pasting}
\begin{proof}
Policies exist, for example a constant action at every date.
Their integrated returns are bounded in absolute value by
\(R/(1-\beta)\). Put \(v=\sup_\pi\int I^\pi\,d\rho\).
For each positive integer \(n\), choose \(\pi^{(n)}\) with
\(\int I^{\pi^{(n)}}d\rho>v-1/n\), and set
\(u=\sup_n I^{\pi^{(n)}}\in M(S)\).
For a positive tolerance \(\delta\), let \(E_n\) consist of
the states where \(n\) is the first index satisfying
\(I^{\pi^{(n)}}\geq u-\delta\). These are Borel and cover
\(S\). Define the decision kernel at a history with first
state in \(E_n\) to be the corresponding kernel of
\(\pi^{(n)}\). Countable pasting preserves all kernel
properties. Starting at that state every finite law and
hence the full law agrees with \(\pi^{(n)}\).
Thus the pasted value is at least \(u-\delta\).

Fix any competing policy \(\pi\). For each \(\delta>0\),
apply the same construction to the countable list consisting
of \(\pi\) and all \(\pi^{(n)}\). Its pasted value \(w_\delta\)
satisfies \(w_\delta\geq\max(u,I^\pi)-\delta\), while
\(\int w_\delta\,d\rho\leq v\).
Also \(v\leq\int u\,d\rho\), since all approaching integrated
values are at most \(\int u\,d\rho\). Consequently
\[
 0\leq\int (I^\pi-u)^+\,d\rho
     =\int\max(u,I^\pi)d\rho-\int u\,d\rho\leq\delta.
\]
Let \(\delta\downarrow0\). A nonnegative function of integral
zero is zero almost everywhere, so \(I^\pi\leq u\)
\(\rho\)-almost everywhere. The pasting with
\(\delta=\varepsilon\) is therefore
\((\rho,\varepsilon)\)-optimal. The conclusion was proved
for each \(\pi\); no uncountable union of exceptional null
sets was taken.
\end{proof}
\end{annexitem}

\begin{annexitem}{annex:p17.blackwell.prior.markov}{thm:p17-blackwell-prior-markov}{Backward replacement with conditional action laws}
\begin{proof}
Choose \(N\geq1\) so that
\(2R\beta^N/(1-\beta)<\varepsilon/2\), except that when
\(R=0\) or \(\beta=0\) any \(N\geq1\) works.
Replace all decisions after \(N\) by a fixed constant action.
The new policy \(\pi^0\) is Markov after \(N\) and satisfies
\(I^{\pi^0}\geq I^\pi-\varepsilon/2\) at every state.
Use \(\eta=\varepsilon/(2N)\) for the following \(N\) replacements.

Suppose the decisions after day \(i\) have already become
deterministic Markov maps. Their restarted return
\(V_{i+1}(s')\) is bounded Borel by the policy-law theorem.
The day-\(i\) contribution to the date-one return is
\[
 w_i(s,a)=\beta^{i-1}\int_S
          [r(s,a,s')+\beta V_{i+1}(s')]q(ds'\mid s,a).
\]
This is bounded Borel. At \(\beta=0\), define the power
at \(i=1\) as \(1\); for \(i>1\) the contribution is zero.
Consider the current joint law of
\((S_1,A_1,\ldots,S_i,A_i)\) from initial law \(\rho\).
Factor its \((S_i,A_i)\)-marginal to obtain a conditional
kernel \(\kappa_i(da\mid s)\), defined at every state.
The selection lemma gives a Borel \(f_i\) such that
\[
 \kappa_i(\{a:w_i(s,a)\geq w_i(s,f_i(s))+\eta\}\mid s)=0
                                                  \quad(s\in S).
\]
Integrating this identity against the current \(S_i\)-law
and using the joint factorization proves
\[
             w_i(S_i,A_i)\leq w_i(S_i,f_i(S_i))+\eta
                        \quad\text{almost surely}.
\]
Condition on \(S_1\). The histories before day \(i\) have
the same law when only its action kernel is replaced.
Their earlier reward is unchanged. The conditional
expectations of the two displayed day-\(i\) functions are
exactly the two future contributions, by the iterated
kernel formula and the Markov form of the later decisions.
Conditional order therefore gives
\(I^{\rm new}(S_1)\geq I^{\rm old}(S_1)-\eta\) almost surely,
which is \(I^{\rm new}\geq I^{\rm old}-\eta\)
outside a \(\rho\)-null set.

Proceed from \(i=N\) to \(1\). Only finitely many null
sets are combined, and the total loss is at most \(N\eta\).
Together with truncation this proves the claimed
\(\varepsilon\)-domination. The selected maps can depend
on \(\rho\), through each current joint law.
For the final optimality assertion, start with a
\((\rho,\varepsilon/2)\)-optimal policy from Theorem 1,
then dominate it within \(\varepsilon/2\) by this construction.
For each competing policy, the union of its first exceptional
set with the domination exceptional set is still null.
\end{proof}
\end{annexitem}

\begin{annexitem}{annex:p17.blackwell.projection.boundaries}{thm:p17-blackwell-projection-boundaries}{A non-Borel projection and all three policy counterexamples}
\begin{proof}
\paragraph{A complete construction of the non-Borel projection.}
Write \(\mathcal N=\mathbb N^{\mathbb N}\), with its product
topology. The metric \(2^{-j}\), where \(j\) is the first
different coordinate, is complete: a Cauchy sequence eventually
fixes each finite prefix, and the resulting sequence is its
limit. The eventually zero sequences are a countable dense
set, so \(\mathcal N\) is Polish.
For a Polish \(X\), call a subset analytic in this proof
if it is the projection of a closed subset of \(X\times\mathcal N\).
We need, and now prove, that every Borel set is analytic.

Closed sets are analytic by taking their product with
\(\mathcal N\). Countable unions of analytic sets are analytic:
if \(F_n\subseteq X\times\mathcal N\) are closed witnesses,
encode the selected \(n\) as the first coordinate of a new
witness and use its remaining coordinates for the witness
in \(F_n\). The resulting set is closed, since a convergent
sequence of witnesses has eventually constant first coordinate
and the corresponding \(F_n\) is closed.
Countable intersections are analytic too: a witness consists
of one sequence for each \(n\), encoded in one sequence by
a fixed bijection \(\mathbb N^2\to\mathbb N\).
The set requiring membership in every \(F_n\) is a countable
intersection of closed inverse images, hence closed.
Its projection is exactly the intersection, using countable
choice to gather witnesses in the reverse inclusion.
Every open subset \(O\) of a metric space is a countable
union of closed sets
\(\{x:d(x,X\setminus O)\geq1/n\}\); treat \(O=X\)
separately. Thus open sets are analytic.
The class of sets whose set and complement are both analytic
is closed under complement and countable union, using the
union and intersection constructions. It contains every closed
set. Hence it contains the whole Borel sigma algebra, proving
the needed assertion.

Fix an enumeration \(e\) of pairs of finite sequences of
the same length, including the empty pair, by natural numbers.
For \(c,x,y\in\mathcal N\), define the closed relation
\[
 \mathcal F(c,x,y)\quad\Longleftrightarrow\quad
       \forall n\geq0\
       c(e(x|n,y|n))>0.
\]
Each individual condition is clopen: the two finite prefixes
and the indicated coordinate of \(c\) determine it.
For every closed \(F\subseteq\mathcal N^2\), assign the
code \(c\) value \(1\) at a pair of prefixes when its
basic cylinder meets \(F\), and \(0\) otherwise.
Then \(\mathcal F(c,x,y)\) holds exactly on \(F\).
One inclusion is immediate. For the other, if every prefix
cylinder meets \(F\), choose points of \(F\) sharing prefixes
of increasing length with \((x,y)\); they converge to it,
and closedness puts it in \(F\).

Let \(U=\{(c,x):\exists y\ \mathcal F(c,x,y)\}\).
Every analytic subset of \(\mathcal N\) is a section
\(U_c\), by the just-proved universality for closed witnesses.
Put \(E=\{x:(x,x)\in U\}\). This set is the projection
of the closed set
\(B_0=\{(x,y):\mathcal F(x,x,y)\}\).
Its complement cannot be analytic. Otherwise it is \(U_c\)
for some \(c\), and substitution at \(c\) gives
\[
        c\notin E\quad\Longleftrightarrow\quad
        c\in U_c\quad\Longleftrightarrow\quad c\in E,
\]
a contradiction. If \(E\) were Borel, its complement would
be Borel and therefore analytic, already proved impossible.
Thus \(B_0\) is a closed set with non-Borel projection.

By \cref{thm:p15-metric-borel-representation},
\(\mathcal N\) has a Borel isomorphism \(\phi\) onto a
Borel subset \(J\) of \(\mathbb R\). Composing with the
homeomorphism
\(t\mapsto(1+t/(1+|t|))/2\) puts \(J\) inside \((0,1)\).
The product map has measurable inverse on \(J^2\), so the
image \(B=(\phi\times\phi)(B_0)\) is Borel in \(J^2\),
and hence in \([0,1]^2\). Its projection is \(\phi(E)\).
Were that projection Borel in \([0,1]\), its inverse image
under \(\phi\) would make \(E\) Borel in \(\mathcal N\).
This proves the required unit-square construction.

\paragraph{Failure of exact attainment with countably many actions.}
For any policy, the first action \(A_1\) is a positive integer
almost surely, so \(E(1/A_1)>0\). Indeed a positive random
variable with zero integral would vanish almost surely,
by the sets on which it exceeds \(1/n\).
Every later reward is at most \(1\); hence
\[
       I^\pi\leq\frac1{1-\beta}-E(1/A_1)
                         <\frac1{1-\beta}.
\]
The constant action \(a\) has return
\((1-1/a)/(1-\beta)\), approaching this bound.
This also handles \(\beta=0\).

\paragraph{No uniformly approximately optimal Borel policy.}
In the fixed-state unit-square example, let
\(g(s)=\int\mathbf1_B(s,a)\pi_1(da\mid s)\).
This is a Borel function, and \(\{g>0\}\subseteq D\).
If it equalled \(D\), the latter would be Borel.
Choose \(s_0\in D\setminus\{g>0\}\); it has first reward
expectation zero. All remaining rewards are at most one,
so \(I^\pi(s_0)\leq\beta/(1-\beta)\).
Choose \(a_0\) with \((s_0,a_0)\in B\).
The constant Borel action \(a_0\) earns
\(1/(1-\beta)\) at \(s_0\). Thus the advantage is at least
one, precluding every \(\varepsilon<1\).
The competitor need optimize only at this selected state.
The value function itself is
\(\mathbf1_D(s)/(1-\beta)\), and is not Borel.

\paragraph{Remembering a coordinate which the current state loses.}
The tagged state space is a Borel subset of a finite
disjoint union of Polish product spaces.
The deterministic transition described in the statement is
Borel on each tagged part, and the reward is a bounded
Borel indicator. Extend its values to impossible transitions
by zero. The remembering policy is Borel: at histories
beginning with \((0,x,y)\) select \(y\) at every date,
and at the other histories choose a fixed default action.
It receives zero first and one on every later date from
an initial \((0,x,y)\in\{0\}\times B\), so its value is
\(\beta/(1-\beta)\).

For a randomized Markov competitor, the second-day action
law depends only on the current tagged state \((1,x)\).
The Borel function
\[
 g(x)=\int\mathbf1_B(x,a)\pi_2(da\mid(1,x))
\]
has \(\{g>0\}\subseteq D\). Choose \(x_0\in D\) with
\(g(x_0)=0\), and choose \(y_0\) witnessing membership
in \(B\). From \((0,x_0,y_0)\), the competitor has zero
reward on each of the first two dates. Its total value
is at most \(\beta^2/(1-\beta)\).
The remembering policy therefore has advantage at least
\(\beta\). For \(0<\beta<1\) this rules out every
uniform domination tolerance smaller than \(\beta\).
When \(\beta=0\), both returns from the tagged graph
are zero, as asserted.
\end{proof}
\end{annexitem}

\begin{annexitem}{annex:p17.blackwell.generated.operators}{thm:p17-blackwell-generated-operators}{Countable choice, finite horizons and the bounded Borel contraction}
\begin{proof}
The joint kernel integral followed by the Borel substitution
\(a=f(s)\) shows \(T_fu\in M(S)\), with
\(\|T_fu\|_\infty\leq R+\beta\|u\|_\infty\).
The same bound and the countable supremum show \(Uu\in M(S)\).
Nonnegative probabilities give monotonicity, and their
total mass one gives the constant-shift identity.
These two properties imply
\[
      \|Bu-Bv\|_\infty\leq\beta\|u-v\|_\infty
\]
by applying monotonicity to
\(u\leq v+\|u-v\|_\infty\) and its reverse.
The full proof of completeness, existence, uniqueness and
both iteration estimates in
\cref{thm:p17-sequential-contraction} applies to the
closed bounded-measurable-function space \(M(S)\).
That lemma explicitly covers arbitrary underlying sets and
their sigma algebras; no finite-state premise is used in
this application.

For a generated \(g\), on its \(n\)-th partition piece,
\(T_gu=T_{f_n}u\leq Uu\). Conversely, for \(\eta>0\)
take the first \(n\) for which \(T_{f_n}u\geq Uu-\eta\).
The countable supremum guarantees existence at each state;
all comparison sets are Borel. Pasting \(f_n\) gives a
Borel generated \(g\) and the desired inequality.

The iterated kernel law gives the \(N\)-day return of
\((g_1,\ldots,g_N)\), with terminal reward \(u\), as
\(T_{g_1}\cdots T_{g_N}u\).
It also gives \(I^{(g,\pi)}=T_gI^\pi\): condition on
the first next state and integrate the restarted future
law. The parameterized policy-law result justifies the
bounded integral even for a history-dependent restarted
\(\pi\). Taking all \(g_i=f\), the terminal tail bound
shows \(T_f^Nu\to I^{f^\infty}\) uniformly.
Hence this return is precisely the fixed point of \(T_f\).

For generated \(g_i\), repeated monotonicity and
\(T_{g_i}\leq U\) give
\(T_{g_1}\cdots T_{g_N}u\leq U^Nu\).
For the reverse supremum inequality, choose positive
\(\eta_1,\ldots,\eta_N\), and choose generated \(g_i\) with
\[
 T_{g_i}U^{N-i}u\geq U^{N-i+1}u-\eta_i.
\]
Backward induction gives
\[
 T_{g_i}\cdots T_{g_N}u
   \geq U^{N-i+1}u-d_i,\qquad
 d_i=\sum_{j=i}^N\beta^{j-i}\eta_j.
\]
At \(i=N\) this is the chosen inequality. If it holds at
\(i+1\), applying \(T_{g_i}\), monotonicity and the
constant-shift identity loses \(\beta d_{i+1}\);
the chosen inequality loses \(\eta_i\), proving the
displayed recursion. Given \(\delta>0\) and \(N\geq1\),
take every \(\eta_i=\delta/N\), so \(d_1\leq\delta\).
Complete the infinite generated policy by any \(f_1\)
after \(N\), which does not change this finite return.
For \(N=0\), the return and \(U^0u\) both equal \(u\)
without a decision. This proves the exact finite-horizon
supremum and simultaneous approximation.

For the countable-action bridge, the map
\((s,n)\mapsto(s,f_n(s))\) is Borel by its countable
pieces. Hence the displayed \(\widehat q,\widehat r\)
have exactly the model's required measurability and bound.
A Borel integer decision \(i_t(s)\) selects the original
map \(g_t(s)=f_{i_t(s)}(s)\), generated by the Borel
partition \(i_t^{-1}\{n\}\). Conversely a partition
exhibiting a generated \(g_t\) defines its integer
decision map by that piece's index. At each state the
corresponding transition measures and realized reward
functions agree. Induction in the iterated kernel law
therefore gives identical joint distributions of states
and rewards at every finite horizon; the uniform tail
bound gives identical infinite returns. Empty partition
pieces and repeated indices cause no change.
\end{proof}
\end{annexitem}

\begin{annexitem}{annex:p17.blackwell.generated.value}{thm:p17-blackwell-generated-value}{Generated upper value and stationary approximation}
\begin{proof}
For a generated policy \(\pi'=(g_1,g_2,\ldots)\),
the preceding monotonicity calculation gives
\[
             T_{g_1}\cdots T_{g_N}u^*
                            \leq U^Nu^*=u^*.
\]
Its distance from \(I^{\pi'}\) is bounded by
\(\beta^N[R/(1-\beta)+\|u^*\|_\infty]\), by the
terminal tail bound. Let \(N\to\infty\) to get
\(I^{\pi'}\leq u^*\).

Choose a generated \(f\) with
\(T_fu^*\geq u^*-\delta\), where
\(\delta=\varepsilon(1-\beta)>0\).
An induction using monotonicity and constant shifts gives
\[
              T_f^Nu^*\geq
                    u^*-\delta\sum_{j=0}^{N-1}\beta^j.
\]
The first step is the hypothesis. Applying \(T_f\) to
the \(N\)-step inequality scales its constant loss by
\(\beta\), and then loses another \(\delta\).
Taking the uniform limit \(I^{f^\infty}\) proves
\(I^{f^\infty}\geq u^*-\delta/(1-\beta)=u^*-\varepsilon\).
The reasoning applies to every generated \(f\) with the
stated residual inequality, proving the last assertion.
\end{proof}
\end{annexitem}

\begin{annexitem}{annex:p17.blackwell.stationary.prior}{thm:p17-blackwell-stationary-prior}{Stationary approximation and the zero-tolerance case}
\begin{proof}
For the prior assertion, take a
\((\rho,\varepsilon/2)\)-optimal deterministic Markov
sequence \(\pi=(f_n)\) by
\cref{thm:p17-blackwell-prior-markov}.
It is generated by itself. Its generated operator has
fixed point \(u^*\geq I^\pi\), and a generated stationary
\(f\) has \(I^{f^\infty}\geq u^*-\varepsilon/2\).
Combining the inequalities gives, for every competitor,
\(I^{\rm competitor}\leq I^{f^\infty}+\varepsilon\)
outside its \(\rho\)-null exceptional set.

For the uniform assertion let \(\pi^*\) be
\(\varepsilon\)-optimal with \(\varepsilon\geq0\), and put
\(v=I^{\pi^*}\). After a fixed first state and action
\((s,a)\), define the continuation policy
\[
 (\pi^*_{s,a})_n(\cdot\mid s_1,a_1,\ldots,s_n)
   =\pi^*_{n+1}(\cdot\mid s,a,s_1,a_1,\ldots,s_n).
\]
For every fixed \((s,a)\) this is an admissible policy.
Its kernels are jointly Borel in the parameter and the
new history. The parameterized policy-law theorem therefore
makes \(w(s,a,s')=I^{\pi^*_{s,a}}(s')\) jointly bounded
Borel. Uniform \(\varepsilon\)-optimality implies
\[
                     w(s,a,s')\leq v(s')+\varepsilon
                                     \quad(s,a,s').
\]
No positive probability of the fixed prefix is required:
the continuation policy is defined by the kernels there.
The first-step law consequently gives, with
\(h(s,a)=\int[r(s,a,s')+\beta v(s')]q(ds'\mid s,a)\),
\[
            v(s)\leq\int h(s,a)\pi^*_1(da\mid s)
                                               +\beta\varepsilon.
\]
The exact mean part of the selection lemma, applied with
any positive auxiliary selection tolerance, produces a
Borel \(f\) with \(h(s,f(s))\geq\int h(s,a)\pi^*_1(da\mid s)\).
Thus \(T_fv\geq v-\beta\varepsilon\).
The same constant-shift induction yields
\[
          T_f^Nv\geq v-\varepsilon\sum_{j=1}^N\beta^j.
\]
Taking limits gives
\(I^{f^\infty}\geq v-\beta\varepsilon/(1-\beta)\).
Every competing return is at most \(v+\varepsilon\);
its advantage over \(I^{f^\infty}\) is therefore at most
\(\varepsilon+\beta\varepsilon/(1-\beta)
 =\varepsilon/(1-\beta)\).
The proof remains valid for \(\varepsilon=0\), since
the mean improvement was exact. It also remains valid
for \(\beta=0\), when only the first selection matters.
\end{proof}
\end{annexitem}

\begin{annexitem}{annex:p17.blackwell.general.verification}{thm:p17-blackwell-general-verification}{All-history verification, Borel value and uniqueness}
\begin{proof}
Suppose \(T_au\leq u\) for all constant actions \(a\).
At each history ending at state \(s\), average the
pointwise inequality
\(\int[r(s,a,s')+\beta u(s')]q(ds'\mid s,a)\leq u(s)\)
against that history's action kernel.
Iterated integration and induction on \(N\) then give
\[
 E_s^\pi\left[\sum_{n=1}^N\beta^{n-1}r(S_n,A_n,S_{n+1})
                      +\beta^Nu(S_{N+1})\right]\leq u(s).
\]
For the induction step, replace the last state's \(u\)
by its next one-step conditional inequality; the
intermediate terms cancel after multiplication by the
appropriate power of \(\beta\). All functions are bounded
Borel, so every integral is defined. The terminal term
has absolute value at most \(\beta^N\|u\|_\infty\).
Letting \(N\to\infty\) proves \(I^\pi(s)\leq u(s)\)
for all policies and all states.

The source also gives a stationary reduction of this
verification argument. Fix \(s\) and \(\delta>0\).
Theorem 6(b), applied to the initial point mass at \(s\),
supplies a stationary \(f^\infty\) with
\(I^\pi(s)\leq I^{f^\infty}(s)+\delta\) for every
competitor. Since \(T_au\leq u\) for all \(a\), one has
\(T_fu\leq u\). Its monotone iterates decrease to
\(I^{f^\infty}\), giving \(I^{f^\infty}\leq u\).
Thus \(I^\pi(s)\leq u(s)+\delta\); let \(\delta\downarrow0\).
The selected \(f\) may depend on \(s,\delta\). This
recovers the source route with precisely its quantifiers.

Now suppose arbitrarily accurate uniformly optimal policies
exist. The stationary replacement theorem produces stationary
\(f_n^\infty\) which are \(1/n\)-optimal: apply it to an
\((1-\beta)/n\)-optimal policy.
Write \(v_n=I^{f_n^\infty}\).
For each \(m,n,s\), \(v_m(s)\leq v_n(s)+1/n\), and
the reverse inequality has error \(1/m\).
Thus \((v_n)\) is uniformly Cauchy. Its bounded Borel
limit \(V\) satisfies
\[
               v_n\leq V,\qquad V-1/n\leq v_n,\qquad
               I^\pi\leq V\quad\text{for every }\pi.
\]
The first two statements also follow directly by taking
the pointwise supremum of the \(1/n\)-optimal inequalities.
Since each \(v_n\) is itself a policy value, this \(V\)
is exactly the supremum of all returns.

Use the sequence \((f_n)\) to define a generated operator
\(U\) with fixed point \(u^*\).
Generated-value optimality gives \(u^*\geq v_n\)
for every \(n\), hence \(u^*\geq V\).
For every \(\delta>0\) it also gives a generated stationary
policy with value at least \(u^*-\delta\); its value
is at most \(V\), hence \(u^*\leq V+\delta\).
Let \(\delta\downarrow0\). Thus \(UV=V\), and
\(\sup_aT_aV\geq V\).
For any fixed action \(a\),
\[
       T_aV\leq T_av_n+\beta/n
             =I^{(a,f_n^\infty)}+\beta/n
             \leq V+\beta/n .
\]
Taking \(n\to\infty\) proves \(T_aV\leq V\) for
every \(a\), so equality holds in the optimality equation.

If \(I^{\pi^*}\) satisfies that equation, the first
verification argument, with \(u=I^{\pi^*}\), shows
that every competing value is below it. Conversely an
optimal policy supplies all positive approximation
tolerances, so the just-proved equation applies to its
return \(V=I^{\pi^*}\).

Finally let \(u,v\in M(S)\) both solve the equation.
Put \(d=\|u-v\|_\infty\). For every \(s,a\),
\(T_au(s)\leq T_av(s)+\beta d\).
Taking suprema over \(a\) gives \(u(s)\leq v(s)+\beta d\).
Interchanging \(u,v\) gives \(d\leq\beta d\), hence
\(d=0\). Both integrals are of actual bounded Borel
functions; no measurability of the operator on other
inputs is invoked.
\end{proof}
\end{annexitem}

\begin{annexitem}{annex:p17.blackwell.essential.optimum}{thm:p17-blackwell-essential-optimum}{Essential coverings and attained Borel maximizers}
\begin{proof}
For any \(s,n,u\), \(T_{f_n}u(s)\) is the constant-action
value \(T_au(s)\) at \(a=f_n(s)\), so
\(Uu\leq\sup_aT_au\).
Conversely every \(a\) at \(s\) is equivalent to
some \(f_n(s)\). Strong equivalence makes their integrands
and transition measures agree; operator equivalence
states the same equality of operators directly.
In either case \(T_au(s)=T_{f_n}u(s)\leq Uu(s)\).
This proves equality of the operators on \(M(S)\).
In particular the supremum is bounded Borel under this
countable covering hypothesis.

Its fixed point \(u^*\) satisfies \(T_au^*\leq u^*\)
for every \(a\). General verification bounds all randomized
history-policy returns by \(u^*\). Generated stationary
policies approximate \(u^*\) uniformly by Theorem 6(a),
so \(u^*\) is the full value and those policies are
\(\varepsilon\)-optimal. Contraction gives uniqueness.

Under essential finiteness, for \(s\in E_m\), every value
\(T_{f_n}u^*(s)\) equals one of the first \(m\) values.
Thus the supremum is attained. Let
\[
 B_n=\{s:T_{f_n}u^*(s)=u^*(s)\}
       \setminus\bigcup_{j<n}\{s:T_{f_j}u^*(s)=u^*(s)\}.
\]
These are Borel and partition \(S\); set \(f=f_n\)
on \(B_n\). Then \(T_fu^*=u^*\).
The unique fixed point of \(T_f\) is \(I^{f^\infty}\),
so this policy attains \(u^*\) everywhere.
For a countable action set enumerate its elements as
constant maps \(f_n\); finite sets can be repeated in an
infinite list and have the single covering piece \(E_m=S\)
for their finite number \(m\). A countable metric subspace
has every subset Borel, so these choices use the specified
Borel action sigma algebra.

For the equivalence criterion, denote the two expected
immediate rewards by \(\bar r_a,\bar r_b\).
Equality of operators at \(u=0\) gives
\(\bar r_a=\bar r_b\). At \(u=\mathbf1_C\), subtract
that equality to obtain
\(\beta[q(C\mid s,a)-q(C\mid s,b)]=0\) for every
Borel \(C\). If \(\beta>0\), the transition measures
therefore agree. Conversely equal measures and equal
expected rewards give equal operators for every bounded
Borel \(u\). If \(\beta=0\), the operator is just
the expected immediate reward and the transition measure
plays no additional role.
For a concrete failure of necessity of pointwise reward
equality, let the common transition put probability
\(1/2\) on each of two states. Rewards \(0,2\) and
\(1,1\) have the same expectation \(1\), and hence
identical operators, but unequal reward functions.
\end{proof}
\end{annexitem}

\begin{annexitem}{annex:p17.blackwell.general.improvement}{thm:p17-blackwell-general-improvement}{Monotone improvement and measurable switching}
\begin{proof}
Let \(v=I^\pi\). The restarted-policy identity gives
\(T_gv=I^{(g,\pi)}\geq v\). Monotonicity yields
\[
             v\leq T_gv\leq T_g^2v\leq\cdots.
\]
The contraction iterates converge uniformly to
\(I^{g^\infty}\). Passing to the limit proves both
inequalities in the first conclusion.

For the second, the two stationary returns are bounded
Borel functions, so their comparison set is Borel and
the prescribed \(h\) is Borel. Let
\(v=\max(I^{f^\infty},I^{g^\infty})\).
Where the first value is at least the second,
\[
        T_hv=T_fv\geq T_fI^{f^\infty}
                    =I^{f^\infty}=v.
\]
On the complementary set the same argument with \(g\)
gives \(T_hv\geq v\). Iterating \(T_h\) and taking its
uniform limit proves \(I^{h^\infty}\geq v\).
The declared first branch includes equality, so the map
is defined everywhere. These arguments use only the
bounded Borel operator properties proved above.
\end{proof}
\end{annexitem}

\begin{annexitem}{annex:p17.blackwell.borel.examples}{thm:p17-blackwell-borel-examples}{Exact computations on the stated spaces}
\begin{proof}
In the first example the state stays at \(s\), so every
policy has reward \(s-1/A_n\) at day \(n\).
Its value is at most
\(s/(1-\beta)-E_s(1/A_1)<s/(1-\beta)\).
The strict inequality holds for every distribution of
the positive integer \(A_1\), by the positivity argument
in the first projection-boundary example.
The constant action \(n\) has value
\((s-1/n)/(1-\beta)\); taking \(n\to\infty\) proves
the asserted supremum and exact error at every state.
Here \(Uu(s)=s+\beta u(s)\), since
\(\sup_n(s-1/n)=s\).
Induction gives \(u_k=s\sum_{j=0}^{k-1}\beta^j\),
including \(u_0=0\); the geometric formula gives the
displayed expression and its uniform convergence.
At \(\beta=0\), this says \(u_k=s\) for \(k\geq1\).
The missing exact optimizer concerns the unattained
action supremum, despite complete measurability.

In the second example, \(\mathbf1_C\) is a Borel action
map. At every state it selects reward one, and the state
never changes. Its reward series sums to \(1/(1-\beta)\).
Every individual reward under any policy is at most one,
which proves optimality. Choosing, for example,
\(C=\mathbb Q\cap[0,1]\) shows that continuity of the
reward or selector is unnecessary.

In the third example, starting at \(0\), constant action
\(0\) earns on dates \(1,3,5,\ldots\), with value
\(1+\beta^2+\beta^4+\cdots=1/(1-\beta^2)\).
Starting at \(1\), it earns on dates \(2,4,6,\ldots\),
with value \(\beta/(1-\beta^2)\). Interchanging the
constant action interchanges these two coordinates.
For \(0<\beta<1\), the larger value at state \(0\)
comes from action \(0\), and at state \(1\) from action
\(1\). Thus the switching rule matches the current
state at every date and earns \(1/(1-\beta)\).
The pointwise maximum of the original two returns is
\(1/(1-\beta^2)\); their difference is
\(\beta/(1-\beta^2)>0\).
At \(\beta=0\) the two values coincide after switching,
which explains the stated positive-discount restriction
for strictness.
\end{proof}
\end{annexitem}
